\exercisesection{}

\begin{enumerate}
\def\labelenumi{\arabic{enumi})}
\tightlist
\item
  Let \(\mathbf{T}\) be a topological space and \(\mu\) a positive measure on \(\mathbf{T}\) . Denote by \(\mathfrak{S}_{\mu}\) the set of subsets of \(\mathbf{T}\) whose boundary is \(\mu\) -negligible.
\end{enumerate}

\begin{enumerate}
\def\labelenumi{\alph{enumi})}
\item
  Show that \(\mathfrak{S}_{\mu}\) is a clan.
\item
  If T is completely regular, every point of T has a fundamental system of neighborhoods contained in \(S_{\mu}\) (observe that if f is a continuous function on T, zero outside a \(\mu\) -integrable set, then the set of real numbers r such that \(f(r)\) is not \(\mu\) -negligible is countable).
\end{enumerate}

\begin{enumerate}
\def\labelenumi{\arabic{enumi})}
\setcounter{enumi}{1}
\tightlist
\item
  Let \(\mathbf{T}\) be a completely regular space and \(\mathfrak{F}\) a filter on \(\mathcal{M}_{+}^{b}(\mathrm{T})\) that converges tightly to a bounded measure \(\mu\) . A Borel subset \(\mathbf{A}\) of \(\mathbf{T}\) is said to be a convergence set (for \(\mathfrak{F}\) ) if \(\lim_{\lambda,\mathfrak{F}} \lambda_{\mathbf{A}} = \mu_{\mathbf{A}}\) .
\end{enumerate}

\begin{enumerate}
\def\labelenumi{\alph{enumi})}
\item
  If the disjoint sets \(\mathbf{A}_1\) and \(\mathbf{A}_2\) are convergence sets, then so is \(\mathbf{A}_1 \cup \mathbf{A}_2\) .
\item
  Let A be an open or closed set. For A to be a convergence set, it is necessary and sufficient that \(\lim_{\lambda,\mathfrak{F}} \lambda^{\bullet}(A) = \mu^{\bullet}(A)\) . If A is a convergence set, then so is \(T - A\) .
\item
  Let A be an open or closed convergence set in T, and B a convergence set such that T - B is also a convergence set. Then A ∪ B and A ∩ B are convergence sets, as are their complements.
\item
  The clan generated by the open convergence sets consists of convergence sets.
\item
  Let \(\mathbf{A}\) be a subset of \(\mathbf{T}\) whose boundary is \(\mu\) -negligible. Then \(\mathbf{A}\) is a convergence set.
\item
  Suppose that T is locally compact or Polish. Show that every compact subset K of T is contained in a compact convergence set (in case T is locally compact, use Exercise 1 b). If T is Polish, use the same exercise to construct a sequence of finite coverings \(U_{p}\) of K by open subsets of T, with \(\mu\) -negligible boundary and of diameter \(<2^{-p}\) ; show that \(L=\bigcap_{p}\bigcup_{U\in\mathfrak{U}_{p}}\overline{U}\) is compact and conclude that it is a convergence set by b)).
\end{enumerate}


\begin{enumerate}
\def\labelenumi{\alph{enumi})}
\setcounter{enumi}{6}
\tightlist
\item
  Extend the result of f) to the case of a convergent sequence of measures on a complete metric space.
\end{enumerate}

\begin{enumerate}
\def\labelenumi{\arabic{enumi})}
\setcounter{enumi}{2}
\tightlist
\item
  Let \(\mathbf{T}\) be a Polish space and \((\mu_n)\) a sequence of bounded positive measures on \(\mathbf{T}\) converging tightly to a measure \(\mu\) . Denote by \(\mathfrak{C}\) the set of subsets \(\mathbf{A}\) of \(\mathbf{T}\) having the following property: for every \(\varepsilon > 0\) , there exists a compact subset \(\mathbf{K}\) of \(\mathbf{A}\) such that \(\sup_{n} \mu_n(\mathbf{A} - \mathbf{K}) < \varepsilon\) . Let \(\mathfrak{D}\) be the set of subsets \(\mathbf{A}\) of \(\mathbf{T}\) that, together with their complement, belong to \(\mathfrak{C}\) .
\end{enumerate}

\begin{enumerate}
\def\labelenumi{\alph{enumi})}
\item
  If the sets A and A' belong to C, then so do \(A \cup A'\) and \(A \cap A'\) ; deduce from this that D is a clan of subsets of T.
\item
  Every subset A of T whose boundary is \(\mu\) -negligible belongs to \(\mathfrak{D}\) (apply Prokhorov's theorem to the interior of A, and use Exer. 2 e)).
\item
  Every \(\mathbf{A} \in \mathfrak{D}\) is a convergence set for the sequence \((\mu_n)\) (cf.~Exer. 2); conversely, every open or closed convergence set belongs to \(\mathfrak{D}\) .
\item
  Let \(\mathbf{A} \in \mathfrak{D}\) . Show that there exist a sequence of disjoint compact sets \(\mathbf{K}_p\) and a subset \(\mathbf{N}\) of \(\mathbf{T}\) having the following properties: \(\alpha\) ) each \(\mathbf{K}_p\) is a convergence set for the sequence \((\mu_n); \beta)\) \(\mathbf{A} \subset \mathbf{N} \cup \bigcup_{p} \mathbf{K}_p\) and \(\bigcup_{p} \mathbf{K}_p \subset \mathbf{A} \cup \mathbf{N}; \gamma)\) for every \(\varepsilon > 0\) , there exists an open neighborhood \(\mathbf{U}\) of \(\mathbf{N}\) such that \(\sup_{n}\mu_n(\mathbf{U}) < \varepsilon\) (apply Exer. 2 f) and Prokhorov's theorem to a suitable subspace of \(\mathbf{T}\) that is the intersection of a sequence of open sets).
\end{enumerate}

\begin{enumerate}
\def\labelenumi{\arabic{enumi})}
\setcounter{enumi}{3}
\item
  Let \(\mathbf{T}\) be a completely regular space, \(t\) a point of \(\mathbf{T}\) , and \(\mathfrak{U}\) a set of Borel subsets of \(\mathbf{T}\) that generates the filter of neighborhoods of \(t\) . Let \(\mathbf{I}\) be a set equipped with a filter \(\mathfrak{F}\) , and let \((\mu_i)_{i \in \mathbb{I}}\) be a family of bounded positive measures of total mass 1 on \(\mathbf{T}\) . In order that \(\lim_{i, \mathfrak{F}} \mu_i = \varepsilon_t\) , it is necessary and sufficient that \(\lim_{i, \mathfrak{F}} \mu_i(\mathrm{U}) = 1\) for every \(\mathrm{U} \in \mathfrak{U}\) .
\item
  Let E be a real Hilbert space admitting an orthonormal basis \((x_{n})_{n\in \mathbf{N}}\) . Let T be the space E equipped with the weakened topology, and let \((a_{n})_{n\in \mathbf{N}}\) be a sequence of real numbers such that \(0 < a_{n} < 1\) for all \(n\) and \(\lim_{n\to \infty}n^2\log a_n = 0\) . For every \(n\in \mathbf{N}\) , set \(\mu_{n} = \sum_{p\in \mathbf{N}}a_{n}^{p}(1 - a_{n})\cdot \varepsilon_{n\cdot x_{p}}\) . Show that \(\mu_{n}\) is a bounded positive measure on T of total mass 1 for every \(n \in \mathbf{N}\) , that \(\lim_{n\to\infty}\mu_{n}=\varepsilon_{0}\) (tight convergence), and that \(\lim_{n\to\infty}\mu_{n}(\mathrm{K})=0\) for every compact subset \(\mathrm{K}\) of T (apply the criterion of the preceding exercise to the set \(\mathfrak{U}\) of neighborhoods of 0 of the form \(\{x \mid (a \mid x) \leqslant 1\}\), where a runs over E). In particular, the set of elements of the sequence \((\mu_{n})_{n \in \mathbf{N}}\) is a relatively compact subset of \(\mathcal{M}_{+}^{b}(\mathbf{T})\) that does not satisfy Prokhorov's condition.
\end{enumerate}


\begin{enumerate}
\def\labelenumi{\arabic{enumi})}
\setcounter{enumi}{5}
\tightlist
\item
  Let \(\mathbf{T}\) be a topological space and \(\mathbf{H}\) a set of positive measures on \(\mathbf{T}\) ; assume that the following condition is fulfilled:
\end{enumerate}

\begin{enumerate}
\def\labelenumi{(\Alph{enumi})}
\setcounter{enumi}{21}
\tightlist
\item
  Every point of \(\mathbf{T}\) admits an open neighborhood \(\mathbf{W}\) such that \(\sup_{\mu \in \mathbf{H}} \mu^{\bullet}(\mathbf{W})\) is finite.
\end{enumerate}

Finally, let \(\mathfrak{U}\) be an ultrafilter on H.

\begin{enumerate}
\def\labelenumi{\alph{enumi})}
\item
  Let \(\mathbf{K}\) be a compact subset of \(\mathbf{T}\) . Show that the induced measures \(\mu_{\mathbf{K}} (\mu \in \mathbf{H})\) converge vaguely with respect to \(\mathfrak{U}\) to a measure \(\pi^{\mathbf{K}}\) on \(\mathbf{K}\) .
\item
  Let \(\mathbf{K}\) and \(\mathbf{L}\) be two compact subsets of \(\mathbf{T}\) such that \(\mathbf{K} \subset \mathbf{L}\) ; show that \((\pi^{\mathbf{L}})_{\mathbf{K}} \geqslant \pi^{\mathbf{K}}\) .
\item
  For every compact subset K of T, let \(i^{\mathbf{K}}\) be the canonical injection of K into T. Show that the family of measures \(i^{\mathbf{K}}(\pi^{\mathbf{K}})\) admits a supremum \(\pi\) in \(\mathcal{M}_{+}(\mathrm{T})\) .
\item
  Let \(f\) be a lower semi-continuous positive function on \(\mathbf{T}\) , and \(g\) an upper semi-continuous positive function on \(\mathbf{T}\) with compact support. Show that \(\pi^{\bullet}(f) \leqslant \lim_{\mu, \mathfrak{U}} \mu^{\bullet}(f)\) and \(\pi^{\bullet}(g) \geqslant \lim_{\mu, \mathfrak{U}} \mu^{\bullet}(g)\) .
\end{enumerate}

¶ 7) Let H be a set of bounded positive measures on a topological space T. Assume that \(\sup_{\mu\in H}\mu^{\bullet}(T)\) is finite, and that for every \(\varepsilon>0\) , there exists a compact subset K of T

such that \(\sup_{\mu \in \mathbf{H}}\mu^{\bullet}(\mathrm{T} - \mathrm{K}) < \varepsilon\) . Show that the condition (V) of the preceding exercise is

verified. Let \(\mathfrak{U}\) be an ultrafilter on H. The measure \(\pi\) is defined as in the preceding exercise.

\begin{enumerate}
\def\labelenumi{\alph{enumi})}
\item
  Show that \(\pi^{\bullet}(g) \geqslant \lim_{\mu, \mathfrak{U}} \mu^{\bullet}(g)\) for every upper semi-continuous bounded positive function \(g\) on T. Deduce from this that \(\pi^{\bullet}(f) = \lim_{\mu, \mathfrak{U}} \mu^{\bullet}(f)\) for every bounded function \(f\) on T whose set of points of discontinuity is \(\pi\) -negligible.
\item
  When T is completely regular, deduce from a) that H is relatively compact for the tight topology (this furnishes a new proof of Th. 1 of No.~5 for the case of positive measures).
\end{enumerate}

¶ 8) Let T be a complete separable metric space, and d its metric. For every closed subset F of T and every real number a > 0, denote by \(F^{a}\) the set of \(x \in T\) such that \(d(x, F) < a\) . Given two bounded positive measures \(\lambda\) and \(\mu\) on T, denote by \(D(\lambda, \mu)\) the infimum of the set of real numbers a > 0 satisfying the inequalities \(\lambda^{\bullet}(F) \leqslant \mu^{\bullet}(F^{a}) + a\) , \(\mu^{\bullet}(F) \leqslant \lambda^{\bullet}(F^{a}) + a\) for every closed subset F of T.

\begin{enumerate}
\def\labelenumi{\alph{enumi})}
\item
  Show that D is a metric on \(\mathcal{M}_+^b\) and that \(\mathrm{D}(\varepsilon_x,\varepsilon_y) = d(x,y)\) for two points \(x\) and \(y\) of T such that \(d(x,y) < 1\) .
\item
  Let \(f\) be a Borel mapping of \(T\) into \(T\) (cf.~TG, IX, §6, No.~3 or GT, IX, §6, Exer. 16) and \(\lambda\) a bounded positive measure on \(T\) . Show that there exists a bounded positive measure \(\mu\) on T such that \(\mu^{\bullet}(\mathrm{A}) = \lambda^{\bullet}\big(f^{-1}(\mathrm{A})\big)\) for every Borel subset A of T. Let \(a > 0\) , and let H be the set of \(x \in \mathrm{T}\) such that \(d(x, f(x)) \geqslant a\) ; show that H is Borel in T and that \(\mathrm{D}(\lambda, \mu) \leqslant \sup (a, \lambda(\mathrm{H}))\) (for every closed subset F of T, one has \(\mathrm{F} \cap (\mathrm{T} - \mathrm{H}) \subset f(\mathrm{F}^a)\) and \(\lambda^{\bullet}(\mathrm{F}) \leqslant \lambda^{\bullet}(\mathrm{F} \cap (\mathrm{T} - \mathrm{H})) + \lambda^{\bullet}(\mathrm{H})\) ).
\item
  Let \(\mu\) and \(\mu_n\) (for \(n \geqslant 1\) ) be bounded positive measures on \(\mathbf{T}\) such that \(\lim_{n \to \infty} \mathrm{D}(\mu_n, \mu) = 0\) . Show that the sequence \((\mu_n)\) converges tightly to \(\mu\) (show that \(\lim_{n \to \infty} \mu_n^\bullet(\mathrm{T}) = \mu^\bullet(\mathrm{T})\) and \(\mu^\bullet(\mathrm{F}) \geqslant \limsup_{n \to \infty} \mu_n^\bullet(\mathrm{F})\) for every closed subset \(\mathbf{F}\) of \(\mathbf{T}\) ; deduce from this that \(\mu^\bullet(f) \leqslant \liminf_{n \to \infty} \mu_n^\bullet(f)\) for every lower semi-continuous function \(f \geqslant 0\) by the method of Lemma 3 of §2, No.~6).
\item
  Conversely, show that for every sequence of bounded positive measures \(\mu_n\) tending tightly to \(\mu\) in \(\mathcal{M}_+^b\) , one has \(\lim_{n \to \infty} D(\mu_n, \mu) = 0\) . One can proceed as follows: let \(\varepsilon > 0\) , and let \(K\) be a compact subset of \(T\) such that \(\mu^\bullet(T - K) < \varepsilon\) and \(\sup_{n} \mu_n^\bullet(T - K) < \varepsilon\) ;
\end{enumerate}

construct a finite family \((\mathrm{B}_{i})_{1\leqslant i\leqslant p}\) of sets with \(\mu\) -negligible boundary, of diameter \(\leqslant\varepsilon/2\) , pairwise disjoint, whose union contains K (cf.~Exer. 1). Let f be a mapping of T into T, constant on each of the sets \(B_{1},\ldots,B_{p}\) and \(\mathbb{C}(B_{1}\cup\cdots\cup B_{p})\) and such that \(f(\mathrm{B}_i) \subset \mathrm{B}_i\) for \(1 \leqslant i \leqslant p\) . Define the measures \(\pi_n\) and \(\pi\) by \(\pi_n^\bullet(\mathrm{A}) = \mu_n^\bullet(\overline{f}(\mathrm{A}))\) and \(\pi^\bullet(\mathrm{A}) = \mu^\bullet(f^{-1}(\mathrm{A}))\) for every Borel subset A of T; deduce from b) the relations \(\mathrm{D}(\pi_n, \mu_n) \leqslant \varepsilon\) , \(\mathrm{D}(\pi, \mu) \leqslant \varepsilon\) and show that \(\lim_{n \to \infty} \mathrm{D}(\pi_n, \pi) = 0\) .

\begin{enumerate}
\def\labelenumi{\alph{enumi})}
\setcounter{enumi}{4}
\tightlist
\item
  The metric D on \(M_{+}^{b}\) is compatible with the tight topology (observe that \(M_{+}^{b}\) is metrizable for the tight topology).
\end{enumerate}

¶ 9) Notations and hypotheses are those of the preceding exercise. Let \((\mu_n)\) be a Cauchy sequence in \(\mathcal{M}_+^b\) for the metric D; show that the sequence \((\mu_n)\) is convergent (this gives a new proof of the fact that the space \(\mathcal{M}_+^b\) is Polish for the tight topology). One can proceed thus:

\begin{enumerate}
\def\labelenumi{\alph{enumi})}
\item
  Let \(\varepsilon > 0\) and \(a > 0\) be two real numbers; there exists a finite subset \(F\) of \(T\) such that \(\sup_{n} \mu_n^\bullet(T - F^a) \leqslant \varepsilon\) (choose an integer \(N \geqslant 1\) and a compact subset \(K\) of \(T\) such that \(\sup_{n \geqslant N} D(\mu_n, \mu_N) \leqslant \varepsilon/2\) and \(\mu_N^\bullet(T - K) \leqslant \varepsilon/2\) ; deduce from this that \(\sup_{n \geqslant N} |\mu_n^\bullet(T) - \mu^\bullet(T)| \leqslant \varepsilon\) and \(\sup_{n \geqslant N} \mu_n^\bullet(T - K^{a/2}) \leqslant \varepsilon\) ; finally, choose a finite subset \(F\) of \(T\) such that \(K^{a/2} \subset F^a\) and \(\sup_{n \geqslant N} \mu_n^\bullet(T - F^a) \leqslant \varepsilon\) ).
\item
  Let \(\varepsilon > 0\) ; there exists a compact subset \(K\) of \(T\) such that \(\sup_{n} \mu_n^\bullet(T - K) \leqslant \varepsilon\) (for every integer \(p \geqslant 1\) choose a finite subset \(F_p\) of \(T\) such that \(\mu_n^\bullet(T - (F_p)^{2^{-p}}) \leqslant \varepsilon / 2^p\) , and set \(K = \bigcap_{p\geqslant 1}(\overline{(F_p)^{2^{-p}}})\) .
\item
  Deduce from \(b\) ) that the Cauchy sequence \((\mu_n)\) has a convergent subsequence (for the metric D).
\end{enumerate}

¶ 10) let T be a completely regular space; denote by S the set of real functions \(f \geqslant 1\) on T, such that the set of points t of T for which \(f(t) \leqslant c\) is compact for every real number c.~For every \(f \in S\) , denote by \(\mathcal{M}_f\) the set of bounded measures \(\mu\) on T such that \(|\mu|^{\bullet}(f) \leqslant 1\) .

\begin{enumerate}
\def\labelenumi{\alph{enumi})}
\item
  For a subset H of \(\mathcal{M}^{b}(T)\) to satisfy Prokhorov's condition, it is necessary and sufficient that there exist an \(f \in S\) such that \(H \subset M_{f}\) (for necessity, take f to be of the form \(c(1 + \sum_{n \geqslant 1} n \cdot f_{n})\) where \(f_{n}\) is the characteristic function of a set \(U_{n}\) such that \(\mathrm{T} - \mathrm{U}_n\) is compact and \(\sup_{\mu \in \mathbf{H}} |\mu|^{\bullet} (\mathrm{U}_n) \leqslant 2^{-n}\) .
\item
  Suppose T is locally compact. For a subset H of \(\mathcal{M}^{b}(T)\) to be relatively compact (for the tight topology), it is necessary and sufficient that there exist a continuous function \(f \in S\) such that \(H \subset M_{f}\) .
\item
  Let C be a closed convex cone in \(\mathcal{M}_{+}^{b}(\mathrm{T})\) . Show that \(\mathrm{C} \cap \mathcal{M}_{f}\) is a cap of C (TVS, II, §7, No.~2, Def. 3) for every \(f \in \mathcal{S}\) , and deduce from this that C is the union of its caps. Denote by E the union of the extremal generators of C. Suppose T is Souslin. Show that for every \(\pi \in \mathrm{C}\) , there exist a Borel subset B of \(\mathcal{M}_{+}^{b}(\mathrm{T})\) contained in E and a positive measure P on B of total mass 1, such that \(\pi = \int_{\mathrm{B}} \mu d\mathrm{P}(\mu)\) (apply Choquet's integral representation theorem).
\end{enumerate}

\begin{enumerate}
\def\labelenumi{\arabic{enumi})}
\setcounter{enumi}{10}
\item
  Let T be a completely regular space. Assume given, for every compact space K and every continuous mapping f of T into K, a positive measure \(\mu_{f,K}\) on K; assume that \(g(\mu_{f,K}) = \mu_{g \circ f,L}\) for any compact spaces K and L and any continuous mappings \(f : T \to K\) and \(g : K \to L\) ; assume moreover that the measure \(\mu_{f,K}\) is concentrated on \(f(T)\) for any f and K. Show that there exists one and only one bounded positive measure \(\pi\) on T such that \(\mu_{f,K} = f(\pi)\) for any f and K (make use of the universal property of the Stone--Čech compactification of T).
\item
  Let T be a completely regular space, a subspace of a locally compact space L. For every measure \(\mu\) (positive or not) on T, there exists a locally compact subspace \(T'\) of L containing T, and a measure \(\mu'\) on \(T'\) concentrated on T that induces \(\mu\) on T.
\end{enumerate}

¶ 13) *Let G be a locally compact abelian group. One denotes by \(\widehat{G}\) the dual group of G and by \(\alpha\) a Haar measure on \(\widehat{G}\) . For every bounded measure \(\mu\) on G, one denotes by \(\mathcal{F}\mu\) the Fourier transform of \(\mu\) , which is a function on the group \(\widehat{G}\) (cf.~Theor. spec., Ch. II, §1, No.~2).

\begin{enumerate}
\def\labelenumi{\alph{enumi})}
\item
  The Fourier transformation \(\mathcal{F}\) is a continuous mapping of \(\mathcal{M}_{+}^{b}(\mathrm{G})\) equipped with the tight topology into \(\mathcal{C}^b (\widehat{\mathrm{G}})\) equipped with the topology of compact convergence. (Make use of the Cor. of Prop. 13 of No.~6.)
\item
  Let \(\mathfrak{F}\) be a filter on \(\mathcal{M}_+^b (\mathrm{G})\) and \(\Phi\) a bounded continuous function on \(\widehat{\mathrm{G}}\) ; assume that \(\lim_{\lambda ,\mathfrak{F}}(\mathcal{F}\lambda)\cdot \alpha = \Phi \cdot \alpha\) (vague convergence) and \(\lim_{\lambda ,\mathfrak{F}}(\mathcal{F}\lambda)(0) = \Phi (0)\) . Show that the filter \(\mathfrak{F}\) converges tightly to a measure \(\mu\) such that \(\mathcal{F}\mu = \Phi\) (show by application of the Stone-Weierstrass theorem that the set of Fourier transforms of continuous functions on G with compact support is a dense linear subspace of \(\mathcal{C}^0 (\overline{\mathrm{G}})\) for uniform convergence; next observe that \(\lim_{\lambda ,\mathfrak{F}}\int (\mathcal{F}\lambda)\cdot u d\alpha = \int \Phi u d\alpha\) for every \(u\in \mathrm{E}\) and that the filter \(\mathfrak{F}\) contains a vaguely compact set of bounded measures; deduce from this that the filter \(\mathfrak{F}\) has at most one cluster point, then that it converges tightly).
\item
  Let \(\mathfrak{F}\) be a filter on \(\mathcal{M}_{+}^{b}(\mathrm{G})\) having a countable base. In order that \(\mathfrak{F}\) converge tightly to a bounded positive measure \(\mu\) , it is necessary and sufficient that there exist a continuous function \(\Phi\) on \(\widehat{\mathbf{G}}\) such that \(\mathcal{F}\lambda\) converges pointwise to \(\Phi\) with respect to \(\mathfrak{F}\) , in which case \(\mathcal{F}\mu = \Phi\) ( \(P\) . Lévy's theorem).*
\end{enumerate}

